\documentclass[UTF8,11pt]{ctexart}
\usepackage[a4paper,margin=24mm]{geometry}
\usepackage{amsmath,amssymb,amsthm,mathtools,mathrsfs}
\usepackage[all]{xy}
\usepackage{xurl,hyperref,enumitem}
\usepackage{tikz}
\usetikzlibrary{arrows.meta,cd,decorations.markings,calc,patterns}
\hypersetup{hidelinks,breaklinks=true}
\setlength{\emergencystretch}{3em}
\allowdisplaybreaks
\newtheorem*{theorem}{Theorem}
\newtheorem*{thm}{Theorem}
\newtheorem*{lemma}{Lemma}
\newtheorem*{proposition}{Proposition}
\newtheorem*{prop}{Proposition}
\newtheorem*{corollary}{Corollary}
\newtheorem*{cor}{Corollary}
\newtheorem*{claim}{Claim}
\theoremstyle{definition}
\newtheorem*{definition}{Definition}
\newtheorem*{defn}{Definition}
\newtheorem*{remark}{Remark}
\newtheorem*{example}{Example}
\newcommand{\R}{\mathbb R}
\newcommand{\C}{\mathbb C}
\newcommand{\Z}{\mathbb Z}
\newcommand{\Q}{\mathbb Q}
\newcommand{\N}{\mathbb N}
\newcommand{\F}{\mathbb F}
\newcommand{\D}{\mathbb D}
\newcommand{\bB}{\mathbb B}
\newcommand{\bC}{\mathbb C}
\newcommand{\bR}{\mathbb R}
\newcommand{\bZ}{\mathbb Z}
\newcommand{\bN}{\mathbb N}
\newcommand{\bQ}{\mathbb Q}
\newcommand{\bD}{\mathbb D}
\newcommand{\bF}{\mathbb F}
\newcommand{\CP}{\mathbb{CP}}
\newcommand{\RP}{\mathbb{RP}}
\newcommand{\sO}{\mathscr O}
\newcommand{\sI}{\mathscr I}
\newcommand{\sF}{\mathscr F}
\newcommand{\sG}{\mathscr G}
\newcommand{\sH}{\mathscr H}
\newcommand{\sR}{\mathscr R}
\newcommand{\sM}{\mathscr M}
\newcommand{\sV}{\mathscr V}
\newcommand{\sU}{\mathscr U}
\DeclareMathOperator{\Hom}{Hom}
\DeclareMathOperator{\Ext}{Ext}
\DeclareMathOperator{\Tor}{Tor}
\DeclareMathOperator{\Spec}{Spec}
\DeclareMathOperator{\Proj}{Proj}
\DeclareMathOperator{\supp}{supp}
\DeclareMathOperator{\im}{im}
\DeclareMathOperator{\id}{id}
\DeclareMathOperator{\Tr}{Tr}
\DeclareMathOperator{\tr}{tr}
\DeclareMathOperator{\rank}{rank}
\DeclareMathOperator{\ord}{ord}
\DeclareMathOperator{\dist}{dist}
\DeclareMathOperator{\End}{End}
\DeclareMathOperator{\Aut}{Aut}
\DeclareMathOperator{\Gal}{Gal}
\DeclareMathOperator{\vol}{vol}
\DeclareMathOperator{\Cl}{Cl}
\DeclareMathOperator{\coker}{coker}
\DeclareMathOperator{\codim}{codim}
\DeclareMathOperator{\divergence}{div}
\DeclareMathOperator{\Ric}{Ric}
\DeclareMathOperator{\grad}{grad}
\DeclareMathOperator{\Hess}{Hess}
\newcommand{\abs}[1]{\left\lvert#1\right\rvert}
\newcommand{\sabs}[1]{\lvert#1\rvert}
\newcommand{\babs}[1]{\bigl\lvert#1\bigr\rvert}
\newcommand{\Babs}[1]{\Bigl\lvert#1\Bigr\rvert}
\newcommand{\bbabs}[1]{\biggl\lvert#1\biggr\rvert}
\newcommand{\BBabs}[1]{\Biggl\lvert#1\Biggr\rvert}
\newcommand{\norm}[1]{\left\lVert#1\right\rVert}
\newcommand{\snorm}[1]{\lVert#1\rVert}
\newcommand{\bnorm}[1]{\bigl\lVert#1\bigr\rVert}
\newcommand{\Bnorm}[1]{\Bigl\lVert#1\Bigr\rVert}
\newcommand{\bbnorm}[1]{\biggl\lVert#1\biggr\rVert}
\newcommand{\BBnorm}[1]{\Biggl\lVert#1\Biggr\rVert}
\newcommand{\widebar}[1]{\overline{#1}}
\newcommand{\nicefrac}[2]{\frac{#1}{#2}}
\newcommand{\linnprod}[2]{\langle#1,#2\rangle}
\newcommand{\myindex}[1]{#1}
\newcommand{\glsadd}[1]{}
\newcommand{\avoidbreak}{}
\newcommand{\sourceref}[1]{\textnormal{[原文：\nolinkurl{#1}]}}
\newcommand{\thmref}[1]{\sourceref{#1}}
\newcommand{\lemmaref}[1]{\sourceref{#1}}
\newcommand{\propref}[1]{\sourceref{#1}}
\newcommand{\corref}[1]{\sourceref{#1}}
\newcommand{\defnref}[1]{\sourceref{#1}}
\newcommand{\exampleref}[1]{\sourceref{#1}}
\newcommand{\exerciseref}[1]{\sourceref{#1}}
\newcommand{\figureref}[1]{\sourceref{#1}}
\newcommand{\sectionref}[1]{\sourceref{#1}}
\newcommand{\appendixref}[1]{\sourceref{#1}}
\newcommand{\stacksref}[2]{\href{https://stacks.math.columbia.edu/tag/#1}{#2 [#1]}}


\newcommand{\E}{\mathbb E}
\newcommand{\Prob}{\mathbb P}
\newcommand{\Reff}{\mathcal R_{\mathrm{eff}}}
\newcommand{\eps}{\varepsilon}
\DeclareMathOperator{\diam}{diam}
\DeclareMathOperator{\Var}{Var}
\DeclareMathOperator{\Crit}{Crit}
\newcommand{\dd}{\,\mathrm d}


\begin{document}
\section*{088\quad 局部紧群Haar测度的覆盖泛函构造与唯一性}
\subsection*{来源与整理范围}
Terence Tao，254A, Notes 3: Haar measure and the Peter-Weyl theorem，2011-09-27；Section 1，Theorem 3 的唯一性与存在性完整证明，选择作者给出的 Tychonoff 紧致性路线。
\par\url{https://terrytao.wordpress.com/2011/09/27/254a-notes-3-haar-measure-and-the-peter-weyl-theorem/}
\par\textbf{恢复方式（整理者说明）：}依据人类原文重新整理以补齐缺失题号；并非旧题证文件的逐字节恢复；2026-09-28。
\subsection*{作者命题与人类证明}

\begin{theorem}[Existence and uniqueness of Haar measure; Theorem 3]
Let $G$ be a $\sigma$-compact locally compact Hausdorff group. There exists a
nonzero left-invariant Radon measure on $G$, unique up to multiplication by
a positive scalar. The analogous assertion holds for right-invariant
measures.
\end{theorem}
\paragraph{Measure convention and representation theorem.}
Here a Radon measure is a Borel measure finite on compact sets, inner
regular and outer regular. We use the Riesz representation theorem:
a positive linear functional on the real space $C_c(G;\R)$ is integration
against a unique Radon measure. Write $I_\mu(f)=\int_G f\,d\mu$ and
$\tau(y)f(x)=f(y^{-1}x)$.

\begin{proof}[Uniqueness]
Let $\mu,\nu$ be left Haar measures. It suffices to show
\[
 I_\nu(f)I_\mu(g)=I_\mu(f)I_\nu(g)\qquad(f,g\in C_c(G)).
\]
Fix $f,g$ and $\epsilon>0$. Uniform continuity on their compact supports
gives an identity neighbourhood $U_\epsilon$ such that
$|f(xy)-f(x)|\leq\epsilon$ and $|g(xy)-g(x)|\leq\epsilon$ for all $x\in G$,
$y\in U_\epsilon$. The neighbourhoods may be chosen inside one fixed compact
set. Every nonempty open set has positive Haar measure. Urysohn's lemma
therefore gives a nonnegative $\psi_\epsilon\in C_c(U_\epsilon)$ with
$\int\psi_\epsilon\,d\mu=1$. Thus
\[
 \int_G f(xy)\psi_\epsilon(y)\,d\mu(y)=f(x)+O(\epsilon)
\]
uniformly in $x$, with supports contained in a fixed compact set.
Integrating against $\nu$, then using left invariance and Fubini, gives
\[
\begin{aligned}
 I_\nu(f)+O(\epsilon)
 &=\int_G\int_G f(xy)\psi_\epsilon(y)\,d\mu(y)\,d\nu(x)\\
 &=\int_G\int_G f(y)\psi_\epsilon(x^{-1}y)\,d\mu(y)\,d\nu(x)\\
 &=\int_G f(y)\left(\int_G\psi_\epsilon(x^{-1})\,d\nu(x)\right)d\mu(y)\\
 &=I_\mu(f)\int_G\psi_\epsilon(x^{-1})\,d\nu(x).
\end{aligned}
\]
The same factor occurs for $g$. Eliminating it yields
$I_\nu(f)I_\mu(g)-I_\mu(f)I_\nu(g)=O(\epsilon)$.
Letting $\epsilon\to0$ proves the assertion, and Riesz representation
identifies $\nu$ as a scalar multiple of $\mu$. Nonzero positivity makes
that scalar positive.
\end{proof}

\begin{proof}[Existence]
It suffices to construct a nonnegative functional on $C_c(G)^+$ that is
positively homogeneous, additive and left-invariant, and is nonzero.
Fix a nonzero $f_0\in C_c(G)^+$. We first construct functionals
$I=I_{f_1,\ldots,f_n,\epsilon}$, for finite families $f_i\in C_c(G)^+$,
that satisfy homogeneity, left invariance, $I(f_0)=1$, a bound
$0\leq I(f)\leq K(f)$ independent of the chosen family and $\epsilon$, and
\[
 |I(f_i+f_j)-I(f_i)-I(f_j)|\leq\epsilon.
\]

Choose a small $\delta>0$ and an identity neighbourhood $U$ so that
$|f_i(xy)-f_i(x)|\leq\delta$ for all $x\in G$, $y\in U$, $0\leq i\leq n$.
Choose a nonzero $\psi\in C_c(G)^+$ supported in $U$. For $f\in C_c(G)^+$
define the covering number
\[
 [f:\psi]=\inf\left\{\sum_{k=1}^Kc_k:
 f\leq\sum_{k=1}^Kc_k\tau(y_k)\psi,\quad c_k\geq0,\ y_k\in G\right\}.
\]
Compactness of the support gives $[f:\psi]<\infty$, and
$[f_0:\psi]>0$ since $f_0\ne0$. Put
\[
 I(f)=\frac{[f:\psi]}{[f_0:\psi]}.
\]
This is homogeneous, left-invariant, monotone and normalized at $f_0$.
Furthermore
\[
 [f:\psi]\leq[f:f_0]\,[f_0:\psi],\qquad
 I(f+g)\leq I(f)+I(g).
\]
Hence the uniform bound holds with $K(f)=[f:f_0]$.

It remains to prove approximate additivity from below. Fix $i,j$. Choose a
cover
\[
 f_i(x)+f_j(x)\leq\sum_{k=1}^Kc_k\psi(y_k^{-1}x),\qquad
 \sum_kc_k\leq\left(I(f_i+f_j)+\frac\epsilon2\right)[f_0:\psi].
\]
Split $c_k=c_k'+c_k''$ by
\[
 c_k'=c_k\frac{f_i(y_k)+\delta}{f_i(y_k)+f_j(y_k)+2\delta},\qquad
 c_k''=c_k\frac{f_j(y_k)+\delta}{f_i(y_k)+f_j(y_k)+2\delta}.
\]
Whenever $\psi(y_k^{-1}x)\ne0$, uniform continuity gives
$|f_i(y_k)-f_i(x)|\leq\delta$ and $|f_j(y_k)-f_j(x)|\leq\delta$.
Consequently
\[
 \frac{f_i(y_k)+\delta}{f_i(y_k)+f_j(y_k)+2\delta}
 \geq\frac{f_i(x)}{f_i(x)+f_j(x)+4\delta}.
\]
Therefore
\[
 \sum_kc_k'\psi(y_k^{-1}x)+4\delta
 \geq\frac{f_i(x)(f_i(x)+f_j(x))}{f_i(x)+f_j(x)+4\delta}+4\delta
 \geq f_i(x).
\]
The analogous estimate holds for $f_j$ and the coefficients $c_k''$.
Take a fixed $g\in C_c(G)^+$ equal to $1$ on the supports of all the $f_i$.
The preceding estimates can be written with $4\delta g$ in place of the
constant $4\delta$: on the support this is the same bound, and outside the
support the required left-hand side is zero. Subadditivity then yields
\[
\begin{aligned}
 I(f_i)&\leq\frac{\sum_kc_k'}{[f_0:\psi]}+4\delta I(g),\\
 I(f_j)&\leq\frac{\sum_kc_k''}{[f_0:\psi]}+4\delta I(g).
\end{aligned}
\]
Thus
\[
 I(f_i)+I(f_j)\leq I(f_i+f_j)+\frac\epsilon2+8\delta I(g)
 \leq I(f_i+f_j)+\frac\epsilon2+8\delta K(g).
\]
Choosing $\delta$ small enough proves the desired approximate additivity
simultaneously for the finite family.

To pass to an exact functional, consider the compact product
\[
 \mathcal K=\prod_{f\in C_c(G)^+}[0,K(f)].
\]
For each finite family and $\epsilon>0$, impose the closed conditions of
homogeneity, left invariance, normalization and the specified approximate
additivity conditions. The functionals just constructed show that these
closed sets have the finite intersection property: for finitely many such
requirements take the union of the families and the smallest tolerance.
Tychonoff's theorem gives a point in their common intersection. This
functional is exactly additive, since the tolerances are arbitrarily small.
Extend it by differences to $C_c(G;\R)$. By the Riesz representation theorem
it is integration against a Radon measure $\mu$. Its normalization makes
$\mu$ nonzero, and the left invariance of the functional makes $\mu$
left-invariant. Inversion, $E\mapsto E^{-1}$, gives the right-invariant
version, completing the theorem.
\end{proof}

\subsection*{整理说明与先修范围（非作者正文）}
范围限定为作者的 $\sigma$-紧、局部紧、Hausdorff 群，不附加作者留作练习的一般化，也不声称左右 Haar 测度相同。网页的系数比估计把应为“$\geq$”的符号排成等号，此处依前后不等式修正。将正泛函的 Riesz 表述明确放在实函数空间上；近似单位明确取非负。保留函数覆盖数的定义与近似加性估计，不只交最终紧致性论证。

先修为 Riesz--Markov 表示、Urysohn 引理、Tychonoff 紧致性、Fubini 与紧支撑连续函数的一致连续性。难点是尚无不变测度时用平移覆盖数构造近似积分，再证明加性误差可同时压低，使拓扑紧致性产生真正的不变积分。
\subsection*{可修改的简短标注（非作者正文）}
$c$：局部紧拓扑群的不变测度

$a$：平移覆盖数；正泛函；近似加性；Tychonoff紧致性；Riesz--Markov表示；Fubini交换积分

\texttt{cross\_theory\_status = yes}。先将群的平移与覆盖转成近似线性泛函，用乘积空间的紧致性取得精确积分，再通过表示定理返回群上的测度；见存在性证明的覆盖数和紧致性两段。

全部标注为非穷尽、可修改初稿。
\subsection*{来源许可与编辑归属}
作者 About 页 Copyright etc. 允许带来源引用复用合理篇幅（例如单篇文章）；本题为其指定小节的编辑转录。许可入口：https://terrytao.wordpress.com/about/ 。
ChatGPT 加入题号、中文来源说明、整理范围和初稿标注；编辑说明与人类证明分列。
\end{document}
